\documentclass[11pt]{article}

\usepackage{amsmath,amssymb,amsthm}

\newcommand{\Obj}{\mathrm{Obj}}
\newcommand{\Hom}{\mathrm{Hom}}
\newcommand{\id}{\mathrm{id}}
\newcommand{\op}{\mathrm{op}}
\newcommand{\Cat}{\mathbf{Cat}}

\newcommand{\FS}[1]{\mathbf{FullSub}(#1)} % full subcategory
\newcommand{\Prod}{\times}

\begin{document}

\title{Torsion Clefts and Variations}
\date{}
\maketitle

\paragraph{Ambient data.}
Fix a category $\mathcal{C}$. Write $\Obj(\mathcal{C})$ for objects and $\Hom_{\mathcal{C}}(X,Y)$ for morphisms.

\paragraph{Full subcategories.}
Given a set $I$ and a function $\mathcal{T} : I \to \Obj(\mathcal{C})$ selecting a family of objects, let
$\FS{\mathcal{T}}$ denote the \emph{full subcategory} of $\mathcal{C}$ on the chosen objects.
We write
\[
  i : \FS{\mathcal{T}} \hookrightarrow \mathcal{C}
\]
for its inclusion functor.
Likewise, for $\mathcal{F} : J \to \Obj(\mathcal{C})$ we write $\FS{\mathcal{F}}$ with inclusion
\[
  q_R : \FS{\mathcal{F}} \hookrightarrow \mathcal{C}.
\]

\paragraph{Adjunction notation.}
An adjunction is written $L \dashv R$.
Its \emph{unit} is a natural transformation $\eta : \id \Rightarrow R L$ and its \emph{counit} is
$\varepsilon : L R \Rightarrow \id$.

\section*{Clefts}

\paragraph{Definition (Cleft).}
Let $I,J$ be sets.
A \emph{cleft} on $\mathcal{C}$ indexed by $(I,J)$ is the following data:

\begin{enumerate}
\item A family $\mathcal{T} : I \to \Obj(\mathcal{C})$ and a family $\mathcal{F} : J \to \Obj(\mathcal{C})$.
\item A functor $i_R : \mathcal{C} \to \FS{\mathcal{T}}$ right adjoint to the inclusion
      $i : \FS{\mathcal{T}} \hookrightarrow \mathcal{C}$, i.e.
      \[ i \dashv i_R. \]
\item A functor $q : \mathcal{C} \to \FS{\mathcal{F}}$ left adjoint to the inclusion
      $q_R : \FS{\mathcal{F}} \hookrightarrow \mathcal{C}$, i.e.
      \[ q \dashv q_R. \]
\end{enumerate}

\medskip
In other words, a cleft is a pair of full subcategories $(\FS{\mathcal{T}},\FS{\mathcal{F}})$ such that
$\FS{\mathcal{T}}$ is \emph{reflective} in $\mathcal{C}$ and $\FS{\mathcal{F}}$ is \emph{coreflective} in $\mathcal{C}$.

\paragraph{The combined functor $K$.}
From a cleft we obtain a functor
\[
  K : \mathcal{C} \to \FS{\mathcal{T}} \Prod \FS{\mathcal{F}}
\]
defined by pairing the two components
\[
  K(X) = (i_R(X),\, q(X)),
\]
and similarly on morphisms.

\paragraph{The essential sequence.}
Define two endofunctors of $\mathcal{C}$:
\[
  i \circ i_R : \mathcal{C} \to \mathcal{C},
  \qquad
  q_R \circ q : \mathcal{C} \to \mathcal{C}.
\]

From the adjunction $i \dashv i_R$ we take the counit
\[
  \varepsilon^i : i \circ i_R \Rightarrow \id_{\mathcal{C}}.
\]
From the adjunction $q \dashv q_R$ we take the unit
\[
  \eta^q : \id_{\mathcal{C}} \Rightarrow q_R \circ q.
\]
Their vertical composite is the natural transformation
\[
  \mathsf{essential} \,=\, \eta^q \circ \varepsilon^i : i \circ i_R \Rightarrow q_R \circ q.
\]
Objectwise, for each $X \in \Obj(\mathcal{C})$, this yields the composable morphisms
\[
  i i_R X \xrightarrow{\ \varepsilon^i_X\ } X \xrightarrow{\ \eta^q_X\ } q_R q X.
\]
If $\mathcal{C}$ also has an initial object $0$ and a terminal object $1$, one may extend this to
\[
  0 \to i i_R X \to X \to q_R q X \to 1,
\]
though in what follows we only require the middle three terms.

\section*{Refinements of a Cleft}

The following definitions package a cleft together with additional structure.

\paragraph{Split clefts.}
A \emph{split cleft} is a cleft equipped with an additional left adjoint to the projection onto the
$\FS{\mathcal{F}}$-component.
Concretely, it consists of:

\begin{enumerate}
\item a cleft $\mathsf{cleft}$,
\item a functor $q_L : \FS{\mathcal{F}} \to \mathcal{C}$,
\item an adjunction
  \[
    q_L \dashv (\pi_R \circ K),
  \]
  where $\pi_R : \FS{\mathcal{T}}\Prod\FS{\mathcal{F}} \to \FS{\mathcal{F}}$ is the right projection.
\end{enumerate}
Equivalently, one may ask for $q_L \dashv q$; since $q = \pi_R \circ K$ on objects and morphisms,
we will freely identify these formulations.

\paragraph{Semidirect clefts.}
A \emph{semidirect cleft} packages a cleft together with data showing that the comparison functor $K$
is monadic.
It consists of:

\begin{enumerate}
\item a cleft $\mathsf{cleft}$,
\item a functor $K_L : \FS{\mathcal{T}}\Prod\FS{\mathcal{F}} \to \mathcal{C}$,
\item an adjunction $K_L \dashv K$,
\item a proof that this adjunction is monadic.
\end{enumerate}

Here ``monadic'' may be read in the following concrete way: the canonical comparison functor from
the domain $\FS{\mathcal{T}}\Prod\FS{\mathcal{F}}$ to the Eilenberg--Moore category of the induced monad on $\mathcal{C}$
admits a weak inverse.

\paragraph{Fibrations of algebras (FoA).}
We use the term \emph{fibration of algebras} as a synonym for \emph{semidirect cleft}.

\paragraph{Rectangular clefts.}
A \emph{rectangular cleft} is a cleft whose $K$ is an equivalence of categories.
The record stores an explicit inverse functor and the witnessing natural isomorphisms:

\begin{enumerate}
\item a cleft $\mathsf{cleft}$,
\item a functor $K^{-1} : \FS{\mathcal{T}}\Prod\FS{\mathcal{F}} \to \mathcal{C}$,
\item natural isomorphisms
  \[
    K \circ K^{-1} \cong \id_{\FS{\mathcal{T}}\Prod\FS{\mathcal{F}}},
    \qquad
    K^{-1} \circ K \cong \id_{\mathcal{C}}.
  \]
\end{enumerate}

From this data one obtains an equivalence
$\mathcal{C} \simeq \FS{\mathcal{T}}\Prod\FS{\mathcal{F}}$.

\paragraph{Torsion theories (stub).}
A \emph{torsion theory} is a cleft on a category $\mathcal{C}$ equipped with a zero object,
together with an assertion that the essential sequence is exact.
Formally it consists of:

\begin{enumerate}
\item a cleft $\mathsf{cleft}$,
\item a zero object structure on $\mathcal{C}$,
\item a field \texttt{essential-exact} (left unspecified here: the intended exactness predicate is to be supplied).
\end{enumerate}

\section*{Semidirect Implies Split (under an initial object)}

Assume in addition that the torsion subcategory $\FS{\mathcal{T}}$ has an initial object $\bot_{\mathcal{T}}$.

\paragraph{Construction.}
Let $\mathsf{sdir}$ be a semidirect cleft, so that we have $K_L \dashv K$.
Consider the functor
\[
  L = \mathrm{const}^\ell(\bot_{\mathcal{T}}) : \FS{\mathcal{F}} \to \FS{\mathcal{T}}\Prod\FS{\mathcal{F}},
\]
which sends $X$ to $(\bot_{\mathcal{T}},X)$.
There is a canonical adjunction
\[
  \mathrm{const}^\ell(\bot_{\mathcal{T}}) \dashv \pi_R,
\]
because morphisms out of the initial object are unique.

Define
\[
  q_L := K_L \circ \mathrm{const}^\ell(\bot_{\mathcal{T}}) : \FS{\mathcal{F}} \to \mathcal{C}.
\]
Then by composition of adjunctions one gets
\[
  q_L \dashv (\pi_R \circ K).
\]
This provides the data required for a split cleft.

\section*{Zero Clefts}

\paragraph{Definition (ZeroCleft).}
A \emph{zero cleft} consists of:

\begin{enumerate}
\item a cleft $\mathsf{cleft}$,
\item initial and terminal objects in each full subcategory $\FS{\mathcal{T}}$ and $\FS{\mathcal{F}}$,
\item proofs that the composites $q \circ i$ and $i_R \circ q_R$ are (naturally isomorphic to) constant functors,
      namely
  \[
    q \circ i \cong \mathrm{const}(\bot_{\FS{\mathcal{F}}}),
    \qquad
    i_R \circ q_R \cong \mathrm{const}(\top_{\FS{\mathcal{T}}}).
  \]
\end{enumerate}

This expresses that the ``cross terms'' vanish in the strongest possible way (by being constant at the relevant
zero-like objects of the full subcategories).

\paragraph{Refinements based on a ZeroCleft.}
One can repeat the previous refinement notions with a \emph{zero cleft} as the base datum:

\begin{itemize}
\item \texttt{SplitZeroCleft}.
\item \texttt{SDirZeroCleft}.
\item \texttt{RectangularZeroCleft}.
\end{itemize}

No other changes are made: each refinement is obtained by forgetting down to the underlying cleft and reusing the same additional fields.

\section*{Null Adjunction}

We conclude with a standard ``null'' adjunction built from initial and terminal objects.

\paragraph{Assumptions.}
Let $\mathcal{C}$ and $\mathcal{D}$ be categories.
Assume $\mathcal{C}$ has an initial object $\bot_{\mathcal{C}}$ and a terminal object $\top_{\mathcal{C}}$,
and similarly $\mathcal{D}$ has $\bot_{\mathcal{D}}$ and $\top_{\mathcal{D}}$.

\paragraph{Definition.}
Define constant functors
\[
  L : \mathcal{C} \to \mathcal{D}, \qquad L = \mathrm{const}(\bot_{\mathcal{D}}),
\]
\[
  R : \mathcal{D} \to \mathcal{C}, \qquad R = \mathrm{const}(\top_{\mathcal{C}}).
\]

The unit $\eta : \id \Rightarrow RL$ is the unique arrow to the terminal object of $\mathcal{C}$,
and the counit $\varepsilon : LR \Rightarrow \id$ is the unique arrow from the initial object of $\mathcal{D}$.
The triangle identities hold by uniqueness of morphisms from an initial object and to a terminal object.

\paragraph{Result.}
These data define an adjunction $L \dashv R$.

\end{document}
